% Pregunta 6. Cuentas públicas y política fiscal (bloque V.3-V.4) · peso alto
\section*{Pregunta 6. Cuentas públicas y política fiscal}

\begin{enunciado}
\textbf{6.} Describa la situación de las \textbf{cuentas públicas} y los principales rasgos de la \textbf{política fiscal}, su relación con el ciclo y su combinación con la política monetaria. ¿Considera que existen riesgos apreciables para la sostenibilidad de la deuda pública? ¿Cómo condicionarán las reglas fiscales europeas vigentes la política fiscal de los próximos años?

\smallskip
\textbf{Peso alto} · \tiempo{15} · guía: V.3 (saldos y tono, reglas, dinámica de la deuda, inflación y cuentas públicas), V.4 (\emph{policy mix}) · cuadros 1 y 9; gráficos 13 y 14.
\end{enunciado}

\begin{examen}
\textit{Las cuentas públicas han mejorado mucho desde la pandemia, pero siguen siendo el desequilibrio macroeconómico más claro de España: un déficit estructural del 3~\% del PIB con la economía por encima de su potencial y una deuda en torno al 100~\% que ha bajado gracias al crecimiento nominal más que al esfuerzo fiscal. La política fiscal ha pasado de procíclica a contracíclica, pero con un ajuste tímido, y las nuevas reglas europeas y el envejecimiento obligan a continuarlo.}

\rotulo{a) La situación de las cuentas públicas}
El déficit baja del 9,9~\% del PIB en 2020 al 2,4~\% en 2025, y la Comisión prevé el 2,4 y el 2,0~\% en 2026-2027 (cuadro~9). Con unos intereses del 2,4~\% del PIB, el saldo primario está equilibrado. La deuda, que pasó del 35,7~\% en 2007 a un máximo del 119,3~\% en 2020, está en el 100,5~\%. Pero el saldo total favorece la imagen: con la brecha en el +2,0~\%, el ciclo aporta en torno a 1,2 puntos, y el saldo ajustado por el ciclo es del $-3{,}6~\%$; sin medidas puntuales, el \textbf{saldo estructural} es del $-3{,}2~\%$ del PIB potencial, y el estructural primario, del $-0{,}8~\%$. Es decir, \uline{con la economía por encima de su potencial, las cuentas públicas aún tienen un déficit estructural del 3~\%}, por encima del 1,5~\% que fijan como referencia las reglas europeas. Los ingresos (42,9~\% del PIB) han subido casi 4 puntos desde 2019 (39,0~\%), y los gastos (45,3~\%) siguen 3 puntos por encima de los de antes de la pandemia (42,0~\%): parte del gasto extraordinario de 2020-2022 se ha quedado como gasto estructural.

La \textbf{inflación} ha ayudado a corto plazo, como predice la teoría: las bases tributarias son nominales y crecen con los precios (y, si los tramos del impuesto sobre la renta no se actualizan, la recaudación crece más que la renta), mientras los intereses solo suben a medida que se renueva la deuda (el tipo implícito pasa del 1,9 al 2,5~\% entre 2021 y 2025, pese a que el bono a 10 años llegó al 3,5~\% en 2023) y las pensiones y los salarios públicos se actualizan con retraso. A medio plazo el efecto se invierte: el gasto indexado y los intereses alcanzan a los ingresos.

\rotulo{b) Los rasgos de la política fiscal y su relación con el ciclo}
El tono se mide con la variación del saldo estructural primario, $\Delta b^{pe}$ (gráfico~13):
\begin{itemize}
\item en \textbf{2012-2013} la consolidación fue muy intensa (+2,9 y +2,5 puntos) con la brecha en su mínimo ($-8{,}7$ y $-9{,}6~\%$): una política \textbf{procíclica}, forzada por la pérdida de acceso a la financiación (diferencial de 435~pb);
\item en \textbf{2015-2019} fue expansiva, primero con la brecha negativa (contracíclica) y después con la brecha positiva ($-1{,}3$ puntos en 2019 con una brecha del +3,1~\%): procíclica de nuevo, porque no se aprovechó la expansión para reducir el déficit estructural y la economía llegó a la pandemia sin colchón;
\item en \textbf{2020-2022}, expansiva frente a la pandemia y a la crisis energética;
\item desde \textbf{2023}, contractiva con la brecha positiva (+1,3, +0,4 y +0,3 puntos en 2023-2025, y +0,4 en 2026-2027): por fin \textbf{contracíclica}, aunque \emph{tímida}, porque un déficit estructural del 3~\% que se reduce 3 o 4 décimas al año sigue siendo, por su nivel, una política holgada.
\end{itemize}
Los estabilizadores automáticos funcionan: en 2025 el ciclo mejora el saldo en 1,2 puntos, y en 2020 lo empeoró en 5.

\rotulo{c) La combinación con la política monetaria}
En 2026 la política monetaria pasa de algo expansiva a neutral para la zona del euro, pero sigue siendo expansiva para España (tipo real negativo, pregunta~4), y la fiscal es ligeramente contractiva. Las 2 van en la dirección correcta frente a una inflación alta y una brecha positiva: el \emph{policy mix} es \textbf{coherente}. Pero como la política monetaria es común y en España es más laxa que en la zona, \uline{la política fiscal es el único instrumento nacional para enfriar la demanda}, y debería apretar algo más de lo previsto; las ayudas frente a la energía, si las hay, deberían ser temporales y focalizadas en los hogares vulnerables. No siempre fue así: en 2018-2019 se combinó una política monetaria muy expansiva con una fiscal expansiva y una brecha positiva, y en 2022 las 2 seguían siendo expansivas con una inflación del 8~\%.

\rotulo{d) ¿Hay riesgos para la sostenibilidad de la deuda?}
La dinámica de la deuda es
\[
d_t-d_{t-1}=(g_t-\tau_t)+\frac{i_t-\gamma^N_t}{1+\gamma^N_t}\,d_{t-1}+\mathrm{ADD}_t .
\]
En 2025 el tipo implícito es del 2,5~\% y el crecimiento nominal, del 5,7~\% (gráfico~14): con $i-\gamma^N=-3{,}2$ puntos, el saldo primario que estabiliza la deuda es un \emph{déficit} del 3,1~\% del PIB. Con el saldo primario equilibrado la deuda debería haber bajado unos 3 puntos; bajó 0,9, lo que deja un ajuste entre déficit y deuda de unos 2 puntos (operaciones financieras que el cuadernillo no permite identificar). Entre 2020 y 2025 la deuda baja 18,8 puntos, y lo hace por el \textbf{efecto bola de nieve}, que la reduce en 31 puntos, frente a unos déficits primarios que suman 9: \uline{la deuda la ha bajado el crecimiento nominal, con la inflación incluida, no el ajuste fiscal}.

Por eso hay \textbf{riesgos apreciables a medio plazo}, aunque no inmediatos. El crecimiento nominal se moderará (5,1 y 4,5~\% previstos) y el tipo de la deuda nueva (el bono a 10 años está en el 3,9~\% en septiembre de 2026) ya supera al implícito (2,5~\%), de modo que el coste medio subirá a medida que se renueve la deuda. Si $i$ y $\gamma^N$ convergen hacia el 3,5~\%, harán falta superávits primarios para que la deuda baje. A ello se suma el gasto en pensiones y sanidad del envejecimiento (la tasa de dependencia sube del 24 al 31~\%). Lo mitigan que la deuda está en euros, un diferencial de solo 45~pb y una economía que crece.

\rotulo{e) Las reglas fiscales europeas}
Las reglas reformadas en 2024 mantienen las referencias del 3~\% de déficit y del 60~\% de deuda, pero se aplican con un solo indicador, una \textbf{senda de gasto primario neto} fijada en un plan fiscal-estructural de medio plazo (4 años, ampliables a 7 con reformas e inversiones), y con 2 salvaguardas: con una deuda superior al 90~\% del PIB, esta debe bajar al menos 1 punto al año de media, y el déficit estructural debe converger hacia un margen del 1,5~\%. Para España suponen 2 cosas: la deuda prevista (99,6 y 98,9~\%) baja 0,9 y 0,7 puntos al año, algo menos de lo que pide la salvaguarda, y el déficit estructural tiene que reducirse en torno a 1,7 puntos. Las reglas obligan, por tanto, a una \textbf{consolidación gradual y sostenida}, de algo más de medio punto al año, mejor por el lado del gasto corriente y ampliando bases que recortando inversión, con cláusulas de escape para circunstancias excepcionales.
\end{examen}

\begin{complemento}{6.1 · Los saldos, del total al estructural primario}
\que{Qué se aplica.} Guía, V.3: las 5 variables del saldo, $b^{pe}=b-\varepsilon\cdot\mathrm{OG}-\text{puntuales}+\text{intereses}$ y $\text{esfuerzo}=\Delta b^{pe}$; «cuidado con el signo»; con déficits muy altos, mirar también el nivel \parencite{EFB2021FiscalStance}; los estabilizadores, $b-b^{ca}=\varepsilon\cdot\mathrm{OG}$.

\begin{tablacaja}
\centering\scriptsize\setlength{\tabcolsep}{3.5pt}
\caption{Del saldo total al esfuerzo fiscal (\% del PIB o del PIB potencial)}\label{tab:saldos}
\begin{tabular}{>{\raggedright\arraybackslash}p{3.5cm}*{9}{>{\centering\arraybackslash}p{0.93cm}}}
\toprule
\rowcolor{tceeazulpalido}
 & \textbf{2019} & \textbf{2020} & \textbf{2021} & \textbf{2022} & \textbf{2023} & \textbf{2024} & \textbf{2025} & \cellcolor{tceegrisclaro}\textbf{2026\,P} & \cellcolor{tceegrisclaro}\textbf{2027\,P} \\
\midrule
Saldo total, $b$ & $-3{,}1$ & $-9{,}9$ & $-6{,}7$ & $-4{,}6$ & $-3{,}3$ & $-3{,}2$ & $-2{,}4$ & $-2{,}4$ & $-2{,}0$ \\
Brecha, OG & 3,1 & $-8{,}5$ & $-3{,}3$ & 0,9 & 1,0 & 1,9 & 2,0 & 1,7 & 1,2 \\
Ajustado por el ciclo, $b^{ca}$ & $-4{,}9$ & $-4{,}8$ & $-4{,}7$ & $-5{,}1$ & $-3{,}9$ & $-4{,}3$ & $-3{,}6$ & $-3{,}5$ & $-2{,}7$ \\
\hspace*{0.8em}$\varepsilon$ implícita, $(b-b^{ca})/\mathrm{OG}$ & 0,59 & 0,60 & 0,61 & 0,58 & 0,64 & 0,61 & 0,59 & 0,60 & 0,60 \\
Estructural, $b^{e}$ & $-4{,}5$ & $-4{,}3$ & $-4{,}6$ & $-5{,}1$ & $-3{,}8$ & $-3{,}5$ & $-3{,}2$ & $-2{,}9$ & $-2{,}5$ \\
Estructural primario, $b^{pe}$ & $-2{,}2$ & $-2{,}1$ & $-2{,}5$ & $-2{,}8$ & $-1{,}5$ & $-1{,}0$ & $-0{,}8$ & $-0{,}4$ & 0,0 \\
\rowcolor{tceesalvia}
Esfuerzo, $\Delta b^{pe}$ & $-1{,}3$ & 0,1 & $-0{,}4$ & $-0{,}3$ & 1,3 & 0,4 & 0,3 & 0,4 & 0,4 \\
Tono y ciclo & \tiny exp. procícl. & \tiny neutral* & \tiny exp. contrac. & \tiny exp. procícl. & \tiny contr. contrac. & \tiny contr. contrac. & \tiny contr. contrac. & \tiny contr. contrac. & \tiny contr. contrac. \\
\bottomrule
\end{tabular}
\fuente{Cuadros~1 y~9 del cuadernillo (AMECO, primavera de 2026). Exp.: expansiva; contr.: contractiva; procícl./contrac.: procíclica o contracíclica según el signo de la brecha. (*) En 2020 las medidas de la pandemia se tratan como temporales y no entran en el saldo estructural (guía, V.3, nota): el saldo estructural no recoge el gran estímulo de ese año. Qué hay que ver: $\varepsilon\approx0{,}6$ en todos los años, la semielasticidad de la Comisión; la diferencia entre $b^{ca}$ y $b^{e}$ (0,4 puntos en 2025) son medidas puntuales que aumentan el déficit; $b^{pe}\approx b^e+\text{intereses}$ ($-3{,}2+2{,}4=-0{,}8$).}
\end{tablacaja}

\que{Error típico.} Juzgar el tono por el saldo total (que mejora 0,8 puntos en 2025 sobre todo por el ciclo y por los intereses) o leer al revés el signo: un $b^{pe}$ que pasa de $-1{,}0$ a $-0{,}8$ \emph{sube}: es contractivo (guía, V.3, subrayado).
\end{complemento}

\begin{complemento}{6.2 · La dinámica de la deuda, con números}
\que{Qué se aplica.} Guía, V.3 (ecuación de la dinámica de la deuda, saldo primario estabilizador, «la diferencia entre el tipo y el crecimiento nominal decide») y «La inflación y las cuentas públicas» (a corto plazo mejora la ratio; a medio plazo el efecto se invierte) \parencite{Blanchard2019Deuda}.

\que{Alternativa de redacción} de la (d) en 2 frases: «La deuda ha bajado 19 puntos desde 2020 sobre todo porque el PIB nominal ha crecido mucho más que el tipo de interés, no por el ajuste. Con el tipo marginal (3,9~\%) ya por encima del implícito (2,5~\%) y el crecimiento nominal moderándose, ese viento de cola se agota: sin superávits primarios la deuda dejará de bajar».

\begin{figuracaja}
\caption{España: variación de la deuda pública y sus componentes, 2021-2027}\label{fig:s08}
\grafico[0.72\linewidth]{s08-deuda}
\fuente{Cuadro~9 del cuadernillo; previsiones de 2026-2027 de la Comisión (tipo implícito previsto: intereses entre la deuda del año anterior). El ajuste déficit-deuda es el residuo. Qué hay que ver: la bola de nieve, muy negativa con la inflación de 2021-2023, pierde fuerza año a año; en 2026-2027 la deuda solo bajaría 0,7-0,9 puntos, menos que el punto anual de la salvaguarda.}
\end{figuracaja}

\begin{tablacaja}
\centering\scriptsize
\caption{Por qué bajó la deuda: descomposición de 2021-2025 (puntos del PIB)}\label{tab:deuda}
\begin{tabular}{>{\raggedright\arraybackslash}p{4.6cm}*{6}{>{\centering\arraybackslash}p{1.25cm}}}
\toprule
\rowcolor{tceeazulpalido}
 & \textbf{2021} & \textbf{2022} & \textbf{2023} & \textbf{2024} & \textbf{2025} & \textbf{2021-2025} \\
\midrule
Tipo implícito, $i$ (\%) & 1,9 & 2,2 & 2,4 & 2,5 & 2,5 & \\
Crecimiento nominal, $\gamma^N$ (\%) & 9,4 & 11,4 & 8,8 & 6,8 & 5,7 & \\
Déficit primario & 4,6 & 2,3 & 0,9 & 0,8 & 0,0 & 8,6 \\
Bola de nieve, $\frac{i-\gamma^N}{1+\gamma^N}d_{t-1}$ & $-8{,}1$ & $-9{,}5$ & $-6{,}4$ & $-4{,}3$ & $-3{,}1$ & $-31{,}5$ \\
Ajuste déficit-deuda (residuo) & $-0{,}1$ & 0,8 & 1,5 & $-0{,}4$ & 2,2 & 4,1 \\
\rowcolor{tceesalvia}
Variación de la deuda, $\Delta d$ & $-3{,}6$ & $-6{,}4$ & $-4{,}0$ & $-3{,}9$ & $-0{,}9$ & $-18{,}8$ \\
\bottomrule
\end{tabular}
\fuente{Cálculo propio con el cuadro~9 (\texttt{calculos.txt}). Ejemplo, 2025: $\frac{0{,}0249-0{,}0572}{1{,}0572}\times101{,}4=-3{,}1$; $\mathrm{ADD}=-0{,}9-0{,}0-(-3{,}1)=+2{,}2$.}
\end{tablacaja}

\begin{tablacaja}
\centering\scriptsize
\caption{Saldo primario que estabiliza una deuda del 100~\% del PIB según $i$ y $\gamma^N$}\label{tab:escenarios}
\begin{tabular}{>{\raggedright\arraybackslash}p{6.0cm}*{3}{>{\centering\arraybackslash}p{1.6cm}}}
\toprule
\rowcolor{tceeazulpalido}
\textbf{Escenario} & $\boldsymbol{i}$ & $\boldsymbol{\gamma^N}$ & \textbf{Saldo primario} \\
\midrule
2025 (observado) & 2,5 & 5,7 & $-3{,}0$ \\
2027 (crecimiento previsto, mismo tipo) & 2,5 & 4,5 & $-1{,}9$ \\
El tipo sube hacia el del bono; el crecimiento se modera & 3,5 & 4,0 & $-0{,}5$ \\
$i=\gamma^N$ & 3,5 & 3,5 & 0,0 \\
Escenario adverso & 4,0 & 3,0 & +1,0 \\
\bottomrule
\end{tabular}
\fuente{Cálculo propio: $\tau-g=\frac{i-\gamma^N}{1+\gamma^N}\,d$ con $d=100$ (guía, A.3.5). Qué hay que ver: el margen que da hoy $i<\gamma^N$ se va estrechando; para bajar la deuda 1 punto al año en el escenario central hará falta un superávit primario de algo más de medio punto.}
\end{tablacaja}

\end{complemento}

\begin{complemento}{6.3 · \emph{Policy mix} y reglas fiscales: cómo se razona}
\begin{figuracaja}
\caption{Los 4 cuadrantes del \emph{policy mix} (infografía de la guía)}\label{fig:policy-mix}
\grafico[0.46\linewidth]{policy-mix}
\fuente{Guía del dictamen, V.4. España en 2026: monetaria expansiva para España aunque neutral para la zona, fiscal ligeramente contractiva; con inflación y brecha positiva, el cuadrante deseable es el de «ambas contractivas».}
\end{figuracaja}

\que{Qué se aplica.} Guía, V.4: «el \emph{policy mix} es coherente si cada política contribuye a los objetivos de la otra, no si las 2 tienen el mismo signo»; en una unión monetaria «el país solo controla la fiscal, que tiene que compensar si su ciclo va por delante del de la zona». V.3: las reglas de 2024 \parencite{RegUE2024Preventivo,RegUE2024Correctivo} y la cláusula para el gasto en defensa \parencite{CE2025Defensa}; recomendaciones por desequilibrio (VI.2).

\que{Las salvaguardas, con números.} Deuda: 100,5 $\to$ 99,6 $\to$ 98,9, es decir, $-0{,}9$ y $-0{,}7$ puntos; la salvaguarda pide $-1{,}0$ de media en el periodo de ajuste. Déficit estructural: de $-3{,}2$ a $-1{,}5$ son 1,7 puntos; a 0,4 puntos al año (el esfuerzo previsto), unos 4 años. Como referencia externa, el FMI proponía un ajuste de en torno a 0,5 puntos del PIB al año entre 2026 y 2030 \parencite{FMI2026ArticuloIV}.

\que{Errores típicos.} Decir que España «incumple» las reglas por tener una deuda del 100~\%: el valor de referencia del 60~\% no se exige de golpe, sino una senda. Y recomendar una consolidación brusca: con el choque energético frenando la demanda, mejor gradual y creíble, anunciada pronto (guía, V.3).

\que{Qué añadir con más tiempo.} Las necesidades brutas de financiación, la vida media de la deuda y quién la tiene (no vienen en el cuadernillo; guía V.3); la calidad del ajuste (proteger la inversión pública y la educación; ampliar bases antes que subir tipos). \que{Qué recortar en el examen.} La lista de etapas del tono (basta: procíclica en 2012-2013 y 2019, contracíclica y tímida desde 2023) y la historia del \emph{policy mix} anterior a 2026.
\end{complemento}

\begin{complemento}{6.4 · Contexto y fuentes: lo que el cuadernillo no trae}
\que{Las reglas, en sus textos.} El brazo preventivo (planes fiscales-estructurales de medio plazo y senda de gasto primario neto) está en el Reglamento (UE) 2024/1263 \parencite{RegUE2024Preventivo} y el correctivo (procedimiento de déficit excesivo) en el Reglamento (UE) 2024/1264 \parencite{RegUE2024Correctivo}; los valores del 3 y del 60~\% siguen en el Protocolo n.º 12 del Tratado \parencite{TFUE2012Protocolo12}.

\que{Defensa.} En 2025 la Comisión invitó a activar de forma coordinada la cláusula de escape nacional para el gasto en defensa, con un margen de hasta el 1,5~\% del PIB \parencite{CE2025Defensa} (guía, V.3). Es un riesgo de gasto y, a la vez, un margen dentro de las reglas: por eso va aquí y no en el texto de examen, que no tiene datos de gasto en defensa.

\que{El fin de los fondos europeos.} El Mecanismo de Recuperación y Resiliencia (Reglamento (UE) 2021/241) acaba en 2026: la inversión pública que financiaba tendrá que pagarse con recursos propios o reducirse, una presión adicional sobre el gasto a partir de 2027.

\que{La recomendación del FMI.} Un ajuste de en torno a 0,5 puntos del PIB al año entre 2026 y 2030 \parencite{FMI2026ArticuloIV} (guía, C.2), coherente con lo que da el cálculo del texto (algo más de medio punto al año).

\que{Lo que el cuadernillo no da y conviene decir que falta}: la vida media de la deuda, quién la tiene (residentes, no residentes, el Eurosistema) y las necesidades brutas de financiación (guía, V.3); el Tesoro los publica en sus informes, pero no en una serie descargable (ver \texttt{datos/fuentes-no-encontrado.md}).
\end{complemento}
